\section{Reflections}

%-----------------------------------------------------------------------
\subsection{Individual Reflections}
%-----------------------------------------------------------------------

\subsubsection*{Samah Daniyal}

The most significant technical shift this project produced for me was 
understanding how agentic workflows fundamentally change what a language model 
can reliably do. Before this project, I understood RAG as a retrieval mechanism; 
fetch documents, pass them to a model, generate a response. What I did not 
appreciate was how much the model's behaviour degrades without explicit 
regulation at each stage: what gets retrieved, whether it is sufficient, whether 
the generated output is actually grounded in it. Each component exists because the model, 
left to itself, will fill gaps with plausible-sounding fabrication. That 
understanding now shapes how I think about any LLM-based system.

The hardest challenge was not technical. Working with an external industry 
partner over a full academic year introduced constraints that no amount of 
engineering can resolve: delayed data access, shifting requirements, and a 
dependency on an organisation with its own priorities and timelines. Navigating 
that, managing expectations, adjusting scope without losing the project's core 
purpose, and maintaining momentum within the team, required a different kind of 
problem-solving than debugging code. It is a skill I did not expect to develop 
through a technical FYP, and one I consider more transferable than most of what 
I built.

The skill I am most confident in now is debugging LLM systems at the theory 
level. Early in the project, encountering unexpected model behaviour felt 
opaque, the model produced something wrong and the cause was unclear. By the 
end, I could trace a failure back to its source: whether it was a retrieval 
threshold set too low, a prompt structure that allowed parametric leakage, or a 
language detection edge case the pipeline had not accounted for. That only 
became possible because I understood the underlying concepts well enough to 
reason about them, not just run the code.

\vspace{1em}
\subsubsection*{Rubab Shah}

In this Final Year Project, I worked on the frontend development and the design of the agentic workflow for an AI-powered public chat portal for Alkhidmat Foundation. My role focused on creating a clean, responsive, and user-friendly interface that could serve a diverse audience, including donors, patients, and beneficiaries, while ensuring smooth integration with the backend AI system.

Alongside the frontend, I designed the agentic workflow that governs how the system processes user queries, retrieves relevant information through the RAG pipeline, and generates responses. This required understanding how to structure interactions between the user, the language model, and external data sources, as well as implementing logic for handling different query types and maintaining context.

The learning curve was significant, especially in integrating frontend components with asynchronous AI responses and in shifting from traditional programming to more dynamic, decision-based workflows. However, through continuous iteration and problem-solving, I developed a stronger understanding of full-stack AI systems.

Overall, this experience helped me improve both my technical and analytical skills, while also giving me insight into building practical AI solutions with real-world impact.

\vspace{1em}
\subsubsection*{Yashfeen Zehra Zaidi}

This project was a fantastic learning curve that allowed me to dive deep into the mechanics 
of LLMs and Retrieval-Augmented Generation (RAG) for Al-Khidmat. By implementing Agentic AI, 
I moved beyond simple prompts to create autonomous workflows capable of handling complex 
queries. I integrated confidence calculation scores to ensure accuracy, which triggers 
human routing whenever the AI’s certainty is low, maintaining a seamless bridge between 
automation and manual support. Beyond these technical milestones, the experience truly 
drove home how much soft skills like clear communication, adaptability, and a strong work 
ethic impact the final result.

\vspace{1em}
\subsubsection*{Fatima Dossa}

The entire FYP was a rewarding experience in terms of both understanding best practices in the 
industry as well as from a research point of view. Collaborating with an industry partner 
like Alkhidmat and being in constant conversation with both supervisors offered us a unique 
perspective on how real projects are executed in the industry; how there are hours of project 
requirement meetings before finalising requirements - and how many different stakeholders are 
involved from developers to product leads. The research in terms of lit review for our 
external RAG based chatbot helped us learn what’s happening recently in AI and we did this 
before implementing any feature from deciding how to categorise domain, confidence calculation, 
retrieval strategies, building the knowledge base, or chunking.

%-----------------------------------------------------------------------
\subsection{Team Reflection}
%-----------------------------------------------------------------------

The team comprised four members, each owning a distinct layer of the system. 
Samah led backend architecture and documentation. Rubab owned the frontend and 
the agentic module. Yashfeen managed UX, testing, and model evaluation. Fatima 
handled the data pipeline and image processing. This division was deliberate: 
the five-layer architecture — presentation, application, AI/ML, data, and 
infrastructure — mapped directly onto individual ownership, which reduced 
coordination overhead during parallel development.

The more demanding challenge was integration. Each layer made assumptions about 
the layers adjacent to it, and those assumptions did not always hold. The 
language detection module, for instance, needed to be tightly coupled to both 
the frontend language state and the backend retrieval strategy — a dependency 
that only became apparent when testing cross-lingual queries end-to-end. 
Resolving these integration points required sustained communication across the 
team rather than isolated problem-solving within each layer.

Working across different approaches to work and different tolerances for 
ambiguity over a year-long project is its own form of engineering. The team 
navigated this by maintaining shared ownership of the core evaluation criteria 
— if the RAG pipeline failed a test case, it was a team problem regardless of 
whose code was involved. That framing kept collaboration functional even when 
individual working styles diverged.

%-----------------------------------------------------------------------
\subsection{Process Reflection}
%-----------------------------------------------------------------------

\textbf{What worked well.} The modular architecture was the correct structural 
decision. Because each component — the Router Agent, the Retriever, the 
Evidence Coverage Agent, the Self-RAG Critic — was implemented as an 
independent unit with defined inputs and outputs, individual components could 
be tested, replaced, and debugged without destabilising the rest of the 
pipeline. When the conversation memory buffer failed (RAG10), the fault was 
isolated and corrected without touching retrieval or generation logic. That 
would not have been possible in a monolithic implementation.

\textbf{What could be improved.} End-to-end latency of 15--60 seconds is the 
system's most significant unresolved limitation. The primary contributors are 
sequential Self-RAG validation stages (each requiring a model call), CPU-bound 
local inference for Urdu queries, and post-generation claim verification. These 
were accepted tradeoffs given the deployment constraints — local inference was 
necessary to maintain offline capability — but they represent a clear target 
for future optimisation. Parallelising the Self-RAG validation stages, moving 
inference to a GPU-enabled server, and streaming partial responses to the 
frontend would collectively bring response times into an acceptable range for 
a production chat interface. The embedding cache addresses repeated queries 
effectively, but does not reduce latency for novel queries, which represent 
the harder problem.

%-----------------------------------------------------------------------
\subsection{Plans vs. Achievements}
%-----------------------------------------------------------------------

\textbf{What was achieved.} The core system — the Agentic RAG engine, the 
multilingual pipeline, the role-based frontend, the human escalation 
mechanism, and the admin analytics dashboard — was fully implemented and 
evaluated. The system handles English, Urdu, and Roman Urdu queries, processes 
image inputs through the OCR pipeline, enforces factual grounding through 
claim-level verification, and escalates low-confidence queries to human agents 
via a ticketing system. These constitute the primary deliverables of the 
original project proposal.

\textbf{What was not achieved.} WhatsApp integration was specified in the 
original proposal as the primary deployment channel, on the basis that 
Alkhidmat Foundation's beneficiaries predominantly communicate through 
WhatsApp. During the project, the industry partner withdrew this requirement 
because they were unable to provide access to the WhatsApp Business API — a 
dependency outside the team's control. The system was consequently deployed as 
a standalone web portal rather than a WhatsApp-native interface. This does not 
affect the RAG pipeline or the agentic architecture, but it does limit the 
system's reach to users who access it through a browser rather than through 
the messaging platform already embedded in their daily use.

\textbf{What was added.} The agentic layer was not part of the original 
proposal. The initial design specified a standard RAG pipeline. During the 
literature review phase, the team identified that standard RAG systems produce 
unacceptable hallucination rates for a deployment context involving vulnerable 
users and sensitive queries. The Router Agent, Self-RAG Critic, and Evidence 
Coverage Agent were added as architectural enhancements in response to that 
finding — a scope expansion that increased implementation complexity but was 
necessary to meet the system's core reliability requirement.

\textbf{Design priorities.} The knowledge base was scoped to Healthcare and 
Donation domains as a deliberate MVP decision. Alkhidmat Foundation operates 
across 12+ service domains, and building a reliable knowledge base across all 
of them within the project timeline — particularly given the delays in data 
access — was not feasible without compromising retrieval quality in every 
domain. Depth in two domains was prioritised over shallow coverage across all 
of them. Broader domain expansion is the most immediate opportunity for future 
development.